\section{Literature Review}
\label{sec:lit_review}

The canonical foundation for fusing heterogeneous biomedical data via plausible-paradoxical reasoning was laid by Quellec et al. \cite{quellec2010dezert}, who applied the Dezert--Smarandache Theory to medical case retrieval and demonstrated that belief-mass combination across modalities could outperform classical Bayesian averaging. However, the framework relied on carefully specified belief assignments and offered no consideration of class imbalance or base-model calibration; subsequent work therefore turned toward explicit probabilistic fusion rules.

Fontani et al. \cite{fontani2013dempster} provided a canonical exposition of Dempster's rule for decision fusion, showing that genuine belief combination proceeds via mass normalization with factor $(1-K)^{-1}$. The authors demonstrated stable fusion under low conflict but flagged that high inter-source conflict ($K \to 1$) triggers unstable normalization that magnifies small input instabilities---an insight that foreshadows the present observation that high cohort conflict ($\bar{K} = 0.4381$) can manufacture artifactual fusion performance. Chen et al. \cite{chen2016calibrating} contributed the methodological foundation for calibrating diagnostic classifier scores to a probability-of-disease scale using Platt, semi-parametric, and isotonic techniques, demonstrating that calibration preserves discrimination and is a separable post-processing step---though the interaction with downstream fusion rules was not then characterized.

Liu et al. \cite{liu2017weighted} extended Dempster--Shafer theory to multimodal settings through a weighted fuzzy framework that propagates source-specific reliability weights through Dempster's combination rule. Baltru{\v{s}}aitis et al. \cite{baltrusaitis2018multimodal} formalized the canonical taxonomy of multimodal machine learning that distinguishes early, intermediate, and late fusion, observing that late (decision-level) fusion is conceptually equivalent to ensemble learning and offers incremental modality integration without retraining. Huang et al. \cite{huang2020fusion} clarified that probability-level ensembling is conceptually a form of early fusion, a distinction under-appreciated in subsequent clinical literature. Nazari et al. \cite{nazari2020ct} established that radiomics alone can yield clinically meaningful ccRCC risk stratification using CT texture features.

Schulz et al. \cite{schulz2021multimodal} introduced the first large-scale multimodal deep-learning ccRCC model, combining histopathology, CT/MRI, and whole-exome sequencing on a TCGA + Mainz cohort (C-index 0.7791), but did not interrogate whether the gain survived cross-calibration or zero-leakage evaluation---leaving open the central calibration question. Stahlschmidt et al. \cite{stahlschmidt2021multimodal} and Kline et al. \cite{kline2022multimodal} confirmed in review that intermediate fusion frequently outperforms unimodal and shallow approaches, while noting that sample-size constraints frequently limit cohort alignment---the bottleneck that bounds every subsequent study. Shao et al. \cite{shao2023deep} advanced deep evidential fusion through Dirichlet uncertainty quantification but applied it to synthetic-to-BCI pipelines rather than clinical settings, again leaving the calibration--fusion interaction uncharacterized.

Within the ccRCC application space, Mahootiha et al. \cite{mahootiha2023multimodal} reported C-index 0.84 / AUROC 0.8 by integrating CT imaging and clinical data; Yang et al. \cite{yang2024multimodal} demonstrated multimodal net benefit on a 251-patient three-center cohort; Ma et al. \cite{ma2025genomic} developed a multiphase CT + WSI fusion achieving AUROC 0.8363 for postoperative recurrence; and Zang et al. \cite{zang2025multimodal} built MPRS on 1,648 ccRCC patients across six centers plus TCGA, achieving internal C-index 0.886 and external 0.838 with KEYNOTE-564 reclassification of 83.3\% of low-risk patients. Chi et al. \cite{chi2026multimodal} demonstrated externally validated interpretable prognostic modeling using automated CT body composition features. Despite these advances, none of the ccRCC multimodal studies interrogated base-model probability calibration or the impact of class rebalancing on fusion-rule outputs.

Andersen et al. \cite{andersen2024impact} provided the most rigorous cross-cohort analysis of class-imbalance correction, evaluating SMOTE, RUS, and ROS on 10 clinical datasets spanning 605,842 patients and showing that SMOTE improves discrimination only marginally ($-0.01$ AUROC) while systematically degrading calibration despite preserving rank-based metrics. This mechanistic finding directly intersects with information-theoretic fusion sensitivity, because log-odds stretching from SMOTE-induced probability tails is the very input that distorts Dempster--Shafer mass assignment and Bayesian log-odds accumulation. Imrie et al. \cite{imrie2025automated} confirmed that late-fusion ensembles are practically advantageous for incremental modality integration in clinical settings, while Xiao et al. \cite{xiao2026urofusion} demonstrated that gated product-of-experts fusion can maintain over 90\% predictive performance under substantial modality dropout---UroFusion-X remains the most architecturally ambitious recent multimodal framework in urological oncology, yet it does not mandate pre-fusion per-fold Platt calibration of base-model probabilities. A comprehensive summary of these contributions is presented in Table~\ref{tab:lit_summary}.

\begin{table*}[t]
\centering
\caption{Summary of Literature Review}
\label{tab:lit_summary}
\small
\renewcommand{\arraystretch}{1.2}
\begin{tabularx}{\textwidth}{@{} p{0.18\textwidth} p{0.06\textwidth} L L @{}}
\toprule
\textbf{Authors} & \textbf{Year} & \textbf{Major Findings} & \textbf{Key Limitations} \\
\midrule
Quellec et al. & 2010 & Belief-mass combination outperforms Bayesian averaging & No calibration; requires carefully specified beliefs \\ \addlinespace
Fontani et al. & 2013 & Canonical rule demonstration; stable under low $K$ & High $K$ amplifies instabilities; clinical imbalance unaddressed \\ \addlinespace
Chen et al. & 2016 & Probability calibration without altering discrimination & Calibration--fusion interaction uncharacterized \\ \addlinespace
Liu et al. & 2017 & Source-specific reliability weighting extends Dempster & DST conflict sensitivity; no calibration characterization \\ \addlinespace
Baltru{\v{s}}aitis et al. & 2018 & Late fusion equivalent to ensemble & Empirical evaluation of calibration--fusion interaction absent \\ \addlinespace
Huang et al. & 2020 & Standardized implementation guidelines & Calibration awareness not emphasized \\ \addlinespace
Nazari et al. & 2020 & Texture entropy correlates with survival & Single-modality; no multimodal extension \\ \addlinespace
Schulz et al. & 2021 & C-index 0.7791; multimodal $>$ unimodal & Cross-calibration not applied; small external cohort \\ \addlinespace
Stahlschmidt et al. & 2021 & Intermediate fusion often outperforms unimodal baselines & Sample-size constraints limit multimodal alignment \\ \addlinespace
Kline et al. & 2022 & Intermediate fusion; late fusion is ensemble & Fusion $\times$ calibration not systematically assessed \\ \addlinespace
Shao et al. & 2023 & Dirichlet uncertainty in DST & Base-probability calibration before BBA uncharacterized \\ \addlinespace
Mahootiha et al. & 2023 & C-index 0.84; AUROC 0.8 & Calibration step not applied \\ \addlinespace
Yang et al. & 2024 & Clinical net benefit & No SMOTE or calibration ablation \\ \addlinespace
Ma et al. & 2025 & AUROC 0.8363 for recurrence & Calibration--fusion interaction not interrogated \\ \addlinespace
Zang et al. & 2025 & Internal C-index 0.886; external 0.838 & Calibration--fusion mechanism uncharacterized \\ \addlinespace
Chi et al. & 2026 & Externally validated interpretable model & Single-modality-leaning; pre-fusion calibration unaddressed \\ \addlinespace
Andersen et al. & 2024 & SMOTE marginally improves discrimination; systematically degrades calibration & Does not extend to downstream fusion-rule evaluation \\ \addlinespace
Imrie et al. & 2025 & Late fusion incrementally integrable & Probabilistic calibration not enforced upstream \\ \addlinespace
Xiao et al. & 2026 & Robust under modality dropout & Pre-fusion Platt calibration not mandatory \\
\bottomrule
\end{tabularx}
\end{table*}

\subsection{Research Gap and the Proposed Work}
Three interrelated gaps motivate the present work. \textbf{Calibration blindness:} across the published ccRCC multimodal literature, fusion gain is reported as a headline metric without systematic interrogation of whether base-model probabilities are calibrated, or whether the gain persists under proper calibration. \textbf{SMOTE-induced calibration degradation has not been propagated through fusion modeling:} although Andersen et al. \cite{andersen2024impact} demonstrated degradation in 10 clinical datasets, no multimodal ccRCC study tests whether SMOTE-induced probability stretching interacts with information-theoretic fusion rules. \textbf{Small alignment cohort size limits fusion benefit observability:} only 126 TCGA-KIRC patients have aligned clinical, genomic, and CT data sufficient for true three-way evaluation, which itself constrains statistical power.

In this article, we present the latest contribution of our research program, \textbf{RenoFusion}\footnote{\url{https://github.com/ali-Hamza817/ccRCC-multimodal-calibration-fusion}}, an open-source multimodal decision-level fusion framework that eliminates these prior limitations through a single rigorous pipeline. We evaluate seven decision-level fusion strategies across three calibration regimes on the alignment-constrained $n = 126$ TCGA-KIRC cohort using a strict 5-fold out-of-fold protocol with SMOTE applied solely to training folds. Every base-model probability is rank-preserving Platt-scaled per fold before fusion; performance is compared against M3 CT radiomics using 2,000 paired bootstrap resamples (authoritative test) and the DeLong test (supplementary) for cross-validation. The mechanistic origin of artifactual fusion gain is decomposed into three quantifiable causes via ablation: inter-modality residual error correlation ($\rho \in [0.3842, 0.4619]$), Dempster--Shafer evidence conflict ($\bar{K} = 0.4381$, 54.8\% high-conflict patients), and SMOTE-induced probability stretching ($+0.0438$ uncalibrated AUROC boost). The resulting finding, that high-dimensional CT radiomics alone (M3 AUROC 0.7037, Platt BSS $+0.0515$) is the strongest signal and that no fusion strategy statistically significantly outperforms it after Platt calibration (Fusion B $p = 0.779$, Dempster--Shafer $p = 0.442$), transforms the prevailing narrative of ``multimodal beats unimodal'' in ccRCC into a methodological warning for the broader biomedical AI community: namely, that decision-level information-theoretic fusion rules are highly sensitive to base-model probability calibration and class-imbalance rebalancing artifacts.
